\documentclass[12pt,a4paper,faculty=ea,language=nl]{ugent-doc}

\usepackage{todonotes}

\usepackage{hyperref}
\hypersetup{
    colorlinks=true,
    urlcolor=black,
    linkcolor=black,
    citecolor=black
}

\usepackage{graphicx}
\graphicspath{{figures/}}

\usepackage{float}

\usepackage{amsmath}

\usepackage{amsthm}

\usepackage{amssymb}

\renewcommand{\contentsname}{Inhoud}
\renewcommand{\figurename}{Figuur}

\setlength{\parindent}{0pt}

\thetitle{\color{ugentblue}Profileren van meerdradige applicaties met PICS}

\infoboxa{\bfseries\large Project in het kader van  het Vakoverschrijdend Project, Groep X}

\infoboxb{Promotor: 
\begin{tabular}[t]{l}
    Prof.\ Dr.\ Ir.\ Lieven Eeckhout\\
\end{tabular}
\\[0.5em]
Begeleiders:
\begin{tabular}[t]{l}
    Ge Liu\\
    Jaime Roelandts\\
\end{tabular}
}

\infoboxc{Studenten: 
\begin{tabular}[t]{lll}
    Tobe De Brabander & \href{mailto:tobe.debrabander@ugent.be}{tobe.debrabander@ugent.be}\\
    Jaco Lagrange & \href{mailto:jaco.lagrange@ugent.be}{jaco.lagrange@ugent.be}\\
    Louis Verstraeten & \href{mailto:louis.verstraeten@ugent.be}{louis.verstraeten@ugent.be}
\end{tabular}
}

\infoboxd{Academiejaar: 2025--2026}

\begin{document}
\maketitle

{\hypersetup{hidelinks}\tableofcontents} 
\newpage

\section{Inleiding}
\todo[inline]{Probleemschetsing in brede context; Algemeen profileren, Tijdsproporitioneel, correcte assignatie van cycli (allemaal zeer algmeen later nog herhaling over dus geen probleem)}
\section{Doelstellingen}
\todo[inline]{Uitleggen wat we analyseren, kort uitleggen wat we willen bereiken ten gevolge van het lezen van wat er al is (PICS die al geimplementeerd zijn, en andere profileringstechnieken), wij willen dit op sniper doen, hier dan ev nog extra puntje over sniper of verwijzing naar apendix waar we blad of 2 over sniper praten en in grote lijnen uitleggen dan het bv een event based simulator is dus dat  cycli tellen anders zal verlopen etc (apendix miss beter om te doen want dan kunnen we er eevnoudig opnieuw later naar verwijzen)}
\section{Planning}
\todo[inline]{planning, dus eerst dips dan commit, er tussen al eens benchmarks etc, dit moet via Gantt chartt ding (geen idee ma staat zo in opgave)}
\section{Technische bespreking}
\todo[inline]{technische details (prob weer verwijzingen nodig naar onze sniper apendix of sniper puntje) we kunnen volgens mij hier wel echt diep ingaan op hoe alles werkt, best ook een extra deel waar we dan eerst voor alles uitleggen gewoon voor eens en voor altijd op een gestructureerde en ni overdreven manier hoe ROB enz werkt, waar we opnieuw verwijzen naar sniper omda het daar uiteraard anders werkt, ma dit zal volgens mij samen me testen en resultaten grootste deel zijn dus we mogen uitgebreid zijn.} 
\section{Testen}
\todo[inline]{hier al onze benchmarks gebruiken om korte testen te doen, best vanaf volgende week (week 7) eigenlijk al want het is ook goed om er in te schrijven da we nog op een probleem kwamen ofz. Ook voor zowel dispatch als commit (of te samen, of alle 3) keer een programma waar we expres iets raar doen optimaliseren.}
\section{resultaten}
\todo[inline]{valt samen met testen.}
\section{reflectie}
\todo[inline]{(dit moet) reflectie op de samenwerking etc of we anders hadden kunnen aanpakken...}
\section{Openstaande problemen}
\todo[inline]{Zien we nog wel}
\section{Besluit}
\section{Bronvermelding}

\end{document}
